\scanpage{546}
\mspage{457}

infini tel que $\operatorname{card} M\leq\pi_0$. Alors il existe un $i_0\in I$ tel que pour \unsure{tt} $i\geq i_0$, $\alpha_{i_0,i}$ soit un iso.

Si on applique ce lemme au cas de $\overline{M'}$ et des $\overline{\alpha_{ji}}$, on trouve donc qu'il existe un $i_0\in I$ tel que les $\overline{\alpha_{ji_0}}$ soient des iso. Cela implique que
$F(\xi)\simeq\underset{i}{\varinjlim}\,\underset{E}{\lfloor\!\underline{\deleted{\unsure{N|V}}\unsure{E|N}_i}\!\rfloor}$ est isomorphe à $E_{i_0}$, donc est de cardinal égal à $\operatorname{card} \unsure{E_{i_0}|E)_{i_0}}=\operatorname{card} E\leq\pi_0$, ce qui achève la démonstration.

D'ailleurs, \unsure{l'|le'}\deleted{\unsure{idée|id}}énoncé \unsure{4.8.13|4.5.13} implique \unsure{4.15.8.1|4.16.8.1}, car dans le système inductif des $(M_i=M,\alpha_{ji})$, de $\varinjlim$ égale à $F(\xi)$, \unsure{S'il|s'il} n'existait pas de $i_0$ comme stipulé dans
4.15.8.1, alors on en conclurait \inserted{aisément} que l'on aurait $\operatorname{Card} F(\xi)>\pi_0$, contrairement à \unsure{4.8.13|4.18.13}. Nous allons dégager ceci dans un lemme :
% archive sheet number: 545
